% ==============================================================================
% chapters/ch04_comparative_analysis.tex
% פרק ד': ניתוח השוואתי – ישראל 2026 מול ארה"ב 2008 והסיכונים הגלומים
% ==============================================================================

\chapter{ניתוח השוואתי – ישראל \num{2026} מול ארה"ב \num{2008} והסיכונים הגלומים}
\label{ch:comparative_analysis}

לאחר שניתחנו בפרק ב' את מנגנוני הכשל של משבר הסאב-פריים האמריקאי (\num{2001}--\num{2008}) ובפרק ג' את התפתחות מבצעי המימון הקבלניים בישראל (\num{2022}--\num{2026}), פרק זה מעמיד את שני השווקים למבחן השוואתי שיטתי. בהתאם למתודולוגיית הניתוח ההשוואתי הרב-ממדי בפרק \num{6.2} ו-\num{6.3} בספר \en{"}מהבעיה אל ה-\en{PDF}\en{"} \cite{segal2026prep}, אנו נמנעים מהקבלה שטחית (\en{"}ישראל היא סאב-פריים \num{2.0}\en{"} או מנגד \en{"}בישראל הכל בטוח\en{"}), ומפרקים את ההשוואה לעשרה פרמטרים כמותיים ומוסדיים. הניתוח מראה כי בעוד שבמישור המיקרו-כלכלי של עיצוב החוזה והתנהגות הרוכש קיים דמיון מדאיג למשבר \num{2008}, הרי שבמישור המאקרו-יציבותי והמשפטי ניצבים בישראל ארבעה \en{"}שוברי גלים\en{"} מוסדיים החוסמים התפתחות של קריסה בנקאית מערכתית, אך מנתבים את מלוא הלחץ אל חוליה פגיעה אחרת: **המגזר היזמי הממונף באשראי חוץ-בנקאי ומשקי הבית בבקבוק המסירות**.

\section{מוקדי הדמיון (אזורי הסכנה): אשליית ההון העצמי, הימור הריבית וניגוד העניינים}
\label{sec:similarities_danger}

השוואה מבנית בין חוזי ה-\en{Hybrid ARM} וה-\en{Option ARM} בארה"ב (\num{2004}--\num{2006}) לבין חוזי ה-\num{20/80}, ה-\num{10/90} והלוואות הבלון בישראל (\num{2022}--\num{2026}) מעלה שלושה מנגנוני כשל משותפים המהווים את \en{"}חומר הבעירה\en{"} בשני השווקים:

\begin{enumerate}[label=\textbf{\arabic*.},rightmargin=0.6cm]
  \item \textbf{אשליית הון עצמי נמוך וכניסת רוכשים שוליים (\en{Marginal Buyers}):} בארה"ב אפשרו הלוואות ה-\en{Piggyback 80/20} רכישת בית עם \pct{0}--\pct{5} הון עצמי \cite{mayer2009rise}. בישראל, אף שהוראות בנק ישראל מחייבות הון עצמי של \pct{25}--\pct{50} בעת נטילת המשכנתא הסופית \cite{boi2025stability}, מבצעי ה-\num{10/90} וה-\num{15/85} אפשרו למשקי בית ולמשקיעים להיכנס לעסקה חוזית מחייבת על נכס בשווי \num{2.8}--\num{3.5} מיליון ש"ח באמצעות העמדת מקדמה ראשונית של \num{280}--\num{420} אלף ש"ח בלבד. בכך נוצר מינוף חוזי אפקטיבי של \num{1:7} עד \num{1:10} על ההון העצמי לאורך תקופת הבנייה.
  \item \textbf{הימור ספקולטיבי על ירידת ריבית או עסקת \en{"}פליפ\en{"} לפני המסירה:} כפי שהוכיחו \en{Gorton (2009)} \cite{gorton2009securitized} ו-\en{Mayer et al. (2009)} \cite{mayer2009rise}, הלוואת ה-\en{2/28 ARM} בארה"ב נלקחה מתוך הנחה שהלווה ימחזר את החוב או ימכור את הבית ברווח לפני איפוס הריבית בחודש \num{24}. באופן זהה לחלוטין, פלח משמעותי מרוכשי ה-\num{20/80} וה-\num{10/90} בישראל בשנים \num{2023}--\num{2025} חתם על החוזים תחת שתי הנחות ספקולטיביות: (א) שריבית בנק ישראל תצנח בחזרה ל-\pct{1.5}--\pct{2.0} עד מועד המסירה ב-\num{2026}--\num{2027}; או (ב) שניתן יהיה להמחות את הזכויות (\en{Assignment / Flip}) לקונה אחר ברווח הון לפני מועד התשלום הסופי. כאשר שתי ההנחות מתבדות מול מלאי שיא של \num{85,900} דירות חדשות מתחרות, האופציה הספקולטיבית פוקעת מחוץ לכסף.
  \item \textbf{ניגוד עניינים של יזמים ומשווקים וניפוח מחיר החוזה (\en{Price Distortion}):} עבודתו של \en{Ben-David (2011)} \cite{ben2011inflated} על שוק הסאב-פריים האמריקאי תיעדה כיצד מוכרים העניקו תמריצים כספיים לקונים דלי-הון כדי לנפח מלאכותית את מחיר העסקה המדווח ולחלץ אשראי גבוה יותר מהבנק. בישראל (\num{2022}--\num{2026}), מנגנון זה שוחזר בקנה מידה תעשייתי: שימור \en{"}מחיר המדף\en{"} הנומינלי תוך סבסוד הלוואת הבלון או ויתור על מדד תשומות הבנייה ($\delta_{\text{NPV}} \approx 10\%--12.4\%$) נועד למנוע הפחתת שווי בדוחות השמאות של הליווי הבנקאי הסגור.
\end{enumerate}

\section{מוקדי השוני: ארבעת \en{"}שוברי הגלים\en{"} המוסדיים של ישראל}
\label{sec:differences_shields}

למרות הדמיון המיקרו-כלכלי במבנה החוזים הדחויים, ניתוח מוסדי ומשפטי מעלה ארבעה הבדלים מבניים תהומיים בין ארצות הברית של \num{2008} לבין ישראל של \num{2026}, המוצגים במטריצה ההשוואתית המלאה בטבלה~\ref{tab:us_vs_israel_comparison}.

\begin{table}[htbp]
\caption{מטריצה השוואתית רב-ממדית: ארה"ב במשבר הסאב-פריים (\num{2001}--\num{2008}) מול שוק הדיור והאשראי בישראל (\num{2022}--\num{2026}) \cite{mian2009consequences,keys2010did,ghent2011recourse,boi2025stability,boi2025directive,salelaw1974}}
\label{tab:us_vs_israel_comparison}
\centering
\footnotesize
\renewcommand{\arraystretch}{1.26}
\rowcolors{2}{BrandGray}{white}
\begin{tabularx}{\textwidth}{R{2.3cm} Y Y R{2.1cm}}
\toprule
\rowcolor{TableHeaderBg}
\textbf{ממד השוואה} & \textbf{ארה"ב (\num{2001}--\num{2008} – משבר הסאב-פריים)} & \textbf{ישראל (\num{2022}--\num{2026} – מבצעי קבלן וקיר המסירות)} & \textbf{כיוון ההשפעה על היציבות בישראל} \\
\midrule
\textbf{1. חוזה המימון המרכזי} & משכנתאות \en{Hybrid 2/28 \& 3/27 ARM}, \en{Option ARM (Neg-Am)} והלוואות \en{Interest-Only}. & דחיית תשלום ליניארית (\num{20/80}, \num{10/90}), מוטציית \num{20/50/30} והלוואות בלון קבלן מסובסדות. & \textbf{\color{BrandGold}סיכון דומה} (דחיית הלם ההחזר ל-\num{36} חודשים). \\
\textbf{2. הון עצמי בחתימה ומגבלת \en{LTV}} & הון עצמי אפסי (\pct{0}--\pct{5}); מינוף מצרפי \en{CLTV} של \pct{95}--\pct{100} באמצעות \en{Piggyback 80/20}. & \pct{10}--\pct{20} בחתימה על החוזה, אך תקרת \en{LTV} קשיחה של בנק ישראל במשכנתא: \pct{75} / \pct{70} / \pct{50}. & \textbf{\color{BrandTeal}יתרון לישראל} לבנקים; \textbf{\color{BrandRed}סיכון סגירה} לרוכש במסירה. \\
\textbf{3. חיתום ואימות הכנסה (\en{PTI})} & מעל \pct{50} הלוואות \en{Low/No-Doc} ו-\en{NINJA}; היעדר תקרת יחס החזר להכנסה מחייבת. & אימות הכנסה מלא ותקרת \en{PTI} קשיחה (\pct{40}--\pct{50}), אך בעסקאות ליניאריות החיתום נדחה למסירה. & \textbf{\color{BrandTeal}מגן חזק לבנקים}; חושף רוכשים לדחיית משכנתא במסירה. \\
\textbf{4. משטר משפטי (\en{Recourse})} & \textbf{\en{Non-Recourse}} במדינות מפתח: הלווה מוסר את המפתח לבנק (\en{Put Option}) ופטור מיתרת החוב. & \textbf{\en{Full Personal Recourse}}: אחריות אישית מלאה של הרוכש + פיצוי מוסכם (\pct{10}--\pct{15}) ליזם בביטול חוזה. & \textbf{\color{BrandTeal}מגן קריטי}: מונע נטישה אסטרטגית (\en{Jingle Mail}). \\
\textbf{5. איגוח ומודל החזקת הסיכון} & \textbf{\en{Originate-to-Distribute}}: מכירת ההלוואות ל-\en{RMBS}, \en{CDO} ו-\en{CDS} במינוף \en{Repo} קצר-מועד. & \textbf{\en{Originate-to-Hold}}: המשכנתאות מוחזקות במלואן במאזני הבנקים המסחריים; אפס איגוח סאב-פריים. & \textbf{\color{BrandTeal}מגן מערכתי}: אין הדבקת בנקאות צללים ונגזרים. \\
\textbf{6. ליווי בנייה והגנת הרוכש} & היעדר חוק מכר פדרלי אחיד; קבלנים ספקולטיביים בנו ללא ליווי סגור וערבויות אישיות לרוכש. & \textbf{חוק המכר (\num{1974})}: חשבון ליווי בנקאי סגור (\en{Ring-Fenced}) וערבות חוק מכר בנקאית לכל שקל ששולם. & \textbf{\color{BrandTeal}מגן מוסדי}: מבודד פרויקטים ומונע אובדן מקדמות. \\
\textbf{7. מינוף יזמים ואשראי חוץ-בנקאי} & חברות בנייה מומנו באשראי בנקאי ושוק ההון עם יכולת מחיקת אופציות קרקע. & יזמים קטנים/בינוניים במינוף קרקע של \pct{70}--\pct{85} ואשראי \textbf{מזנין חוץ-בנקאי} בריבית \en{Prime+5\%--8\%}. & \textbf{\color{BrandRed}מוקד פגיעות חריף בישראל} (\en{Tier-2/Tier-3}). \\
\textbf{8. כריות ההון של המערכת הבנקאית} & בנקי השקעות במינוף \num{1:30}--\num{1:40} (הון עצמי מוחשי של \pct{2.5}--\pct{3.5}); בנקים מסחריים עם כריות דקות. & חמישה בנקים ריכוזיים תחת פיקוח הדוק עם יחס הון עצמי רובד \num{1} (\textbf{\en{CET1}}) של \textbf{\pct{10.6}--\pct{10.9}}. & \textbf{\color{BrandTeal}עמידות גבוהה}: כרית ספיגת הפסדים רחבה מאוד. \\
\bottomrule
\end{tabularx}
\end{table}

\subsection{ניתוח מתמטי של אפקט ה-\en{Full Recourse} וחוק המכר בישראל מול ארה"ב}
ההבדל המכריע ביותר בהתנהגות משקי הבית טמון במשטר ה-\en{Full Recourse} בישראל לעומת ה-\en{Non-Recourse} בארה"ב \cite{ghent2011recourse}. כפי שהראינו במשוואה~\eqref{eq:non_recourse_put} בפרק ב', בארה"ב ברגע שמחיר הדירה $P_T$ ירד מתחת ליתרת החוב $L_T$, הלווה מימש את אופציית הנטישה ומסר את המפתח לבנק.

לעומת זאת, בישראל חוזה הרכישה מול הקבלן וחוזה המשכנתא מול הבנק מקימים חבות אישית מלאה (\en{Full Personal Liability}). אם רוכש דירה במבצע \num{20/80} (ששילם מקדמה $\alpha_0 P_{\text{nom}}$ במועד החתימה ונותר חייב $\left(1-\alpha_0\right)P_{\text{nom}}$ במועד המסירה $T$) מבקש להפר את החוזה וליטוש את העסקה עקב ירידת שווי השוק של הדירה ל-$P_T < P_{\text{nom}}$, היזם (והבנק המלווה) רשאים לחלט פיצוי מוסכם קבוע בחוזה בשיעור $\kappa \in [0.10, 0.15]$ ממלוא מחיר החוזה $P_{\text{nom}}$, ובנוסף לתבוע את הפרש הירידה בשווי הנכס מתוך נכסיו האישיים והכנסתו העתידית של הרוכש. לפיכך, תנאי הכדאיות הכלכלית לנטישה אסטרטגית מרצון בישראל דורש:
\begin{equation}
\label{eq:israel_recourse_condition}
\text{Loss}_{\text{complete}}\left(T\right) = P_{\text{nom}} - P_T + \text{PV}\left(\Delta M_{\text{mortgage}}\right) > \kappa P_{\text{nom}} + \theta_{\text{recourse}}\max\left(P_{\text{nom}}\left(1-\kappa\right) - P_T,\, 0\right) + C_{\text{legal}}
\end{equation}
כאשר $\theta_{\text{recourse}} \approx 1$ מייצג את שיעור הגבייה האפקטיבי בהליכי הוצאה לפועל מול משק בית בעל הכנסה או נכסים בישראל. ממשוואה~\eqref{eq:israel_recourse_condition} נובעת תוצאה אנליטית מרכזית: **בישראל אין לרוכש אופציית \en{Put} חינמית על הדירה**. גם אם מחיר הדירה בשוק ירד ב-\pct{10}--\pct{12}, רוכש ישראלי לא \en{"}יזרוק את המפתחות\en{"} מרצונו החופשי, אלא יעשה כל מאמץ למכור את הדירה בשוק המשני (אף בהפסד של רוב ההון העצמי שהפקיד), לגייס הלוואות משפחתיות או לפרוס את המשכנתא ל-\num{30} שנה — וייקלע לחדלות פירעון אך ורק כאשר נחסמת בפניו רגולטורית האפשרות לקבל משכנתא מהבנק (\en{Inability to Close}) וגם נכשלה מכירת הדירה לצד שלישי (\en{Double Liquidity Lock}).

\begin{figure}[htbp]
\centering
\begin{english}
\begin{tikzpicture}
\begin{axis}[
    xbar,
    width=13.8cm,
    height=7.8cm,
    xmin=0, xmax=10.5,
    ytick={1,2,3,4,5,6,7},
    yticklabels={
        {\textbf{7. Bank Capital Thinness (CET1)}},
        {\textbf{6. Securitization Contagion (CDO/Repo)}},
        {\textbf{5. Strategic Default Risk (Non-Recourse)}},
        {\textbf{4. Underwriting Laxity (NINJA/No-Doc)}},
        {\textbf{3. Developer \& Mezzanine Leverage}},
        {\textbf{2. Price-to-Income \& Affordability Stretch}},
        {\textbf{1. Deferred Payment Shock (ARM / 20-80)}}
    },
    xlabel={\textbf{Systemic Vulnerability Score (1 = Minimal Risk, 10 = Extreme Vulnerability)}},
    bar width=7.5pt,
    grid=major,
    grid style={dashed, gray!35},
    legend style={at={(0.5,1.08)}, anchor=south, legend columns=2, font=\small, draw=BrandNavy},
    nodes near coords,
    nodes near coords style={font=\scriptsize\bfseries},
    enlarge y limits=0.12
]
\addplot[fill=BrandRed!80, draw=BrandRed] coordinates {
    (8.8,1) (7.5,2) (6.8,3) (9.5,4) (9.2,5) (9.8,6) (9.0,7)
};
\addlegendentry{US Subprime Crisis (2007--2008)}

\addplot[fill=BrandBlue!80, draw=BrandNavy] coordinates {
    (8.4,1) (8.9,2) (8.6,3) (3.8,4) (2.2,5) (1.2,6) (2.5,7)
};
\addlegendentry{Israel Housing \& Delivery Wall (2026--2027)}
\end{axis}
\end{tikzpicture}
\end{english}
\caption{פרופיל פגיעות השוואתי רב-ממדי (ציון \num{1}--\num{10}): ארה"ב במשבר הסאב-פריים (\num{2007}--\num{2008}) מול ישראל לקראת קיר המסירות (\num{2026}--\num{2027}) בשבעה ממדי סיכון מערכתיים \cite{segal2026prep,boi2025stability}}
\label{fig:comparative_risk_profile}
\end{figure}

\section{מפת הסיכונים הייחודית לישראל: היכן כן עלול להיווצר כשל בקיר המסירות?}
\label{sec:israel_unique_risk_map}

פרופיל הפגיעות ההשוואתי באיור~\ref{fig:comparative_risk_profile} מדגים בבירור את האסימטריה של המקרה הישראלי: בעוד שבארה"ב ציוני הפגיעות היו גבוהים מאוד (\num{8.8}--\num{9.8}) בממדי האיגוח, היעדר החיתום, משטר ה-\en{Non-Recourse} ודקות ההון הבנקאי, בישראל ממדים אלו מוגנים היטב (ציוני פגיעות נמוכים של \num{1.2}--\num{3.8}). לעומת זאת, ישראל מציגה פגיעות גבוהה במיוחד — ואף חמורה יותר מארה"ב — בשלושה מוקדים ספציפיים המרכיבים את **מפת קווי השבר של \num{2026}--\num{2027}** (איור~\ref{fig:israel_fault_line_map}):

\begin{figure}[htbp]
\centering
\begin{english}
\begin{tikzpicture}[
    node distance=0.55cm and 0.6cm,
    faultbox/.style={
        rectangle,
        rounded corners=4pt,
        draw=BrandRed,
        line width=1.1pt,
        fill=BrandLightRed,
        text width=4.4cm,
        minimum height=2.3cm,
        align=center,
        font=\footnotesize
    },
    shieldbox/.style={
        rectangle,
        rounded corners=4pt,
        draw=BrandTeal,
        line width=1.2pt,
        fill=BrandLightTeal,
        text width=14.2cm,
        minimum height=1.2cm,
        align=center,
        font=\small
    },
    corebox/.style={
        rectangle,
        rounded corners=4pt,
        draw=BrandNavy,
        line width=1.2pt,
        fill=BrandLightBlue,
        text width=14.2cm,
        minimum height=1.1cm,
        align=center,
        font=\small\bfseries
    },
    arrow/.style={-{Stealth[length=2.6mm,width=2.0mm]}, line width=1.1pt, color=BrandRed},
    blockarrow/.style={-{Bar[width=4mm,line width=1.5pt]}, line width=1.2pt, color=BrandTeal}
]
    \node[faultbox] (f1) {\textbf{\color{BrandRed}Fault Line A: Buyers}\\[3pt]\texthebrew{רוכשי \num{20/80} ו-\num{10/90} שנדחים במבחן \en{PTI > 40\%} או משפרי דיור שנתקעו עם דירתם הישנה}};
    \node[faultbox, right=of f1] (f2) {\textbf{\color{BrandRed}Fault Line B: Inventory}\\[3pt]\texthebrew{\num{85,900} דירות חדשות לא-מכורות מתחרות המונעות מכירת \en{"}פליפ\en{"} ברווח לפני טופס \num{4}}};
    \node[faultbox, right=of f2] (f3) {\textbf{\color{BrandRed}Fault Line C: Tier-3 Devs}\\[3pt]\texthebrew{יזמים קטנים/בינוניים במינוף קרקע ומזנין חוץ-בנקאי (\en{Prime+6\%}) עם \en{DSCR < 1.0}}};

    \node[shieldbox, below=0.75cm of f2] (shield) {\textbf{\color{BrandTeal}Israeli Institutional Firewall (Sale Law + BOI LTV/PTI + Full Recourse + CET1 10.85\%)}\\[2pt]\texthebrew{בולמת הדבקה מערכתית לבנקים, אך כולאת את ההפסד אצל היזמים הממונפים, קרנות המזנין וההון העצמי של הרוכשים}};

    \node[corebox, below=0.55cm of shield] (bank) {\color{BrandNavy}Commercial Banking Core Stability Preserved (\texthebrew{יציבות המערכת הבנקאית נשמרת: \en{CET1 > 9.38\%} גם בתרחיש קיצון})};

    \draw[arrow] (f1.south) -- ([xshift=-4.8cm]shield.north);
    \draw[arrow] (f2.south) -- (shield.north);
    \draw[arrow] (f3.south) -- ([xshift=4.8cm]shield.north);
    \draw[blockarrow] (shield.south) -- (bank.north);
\end{tikzpicture}
\end{english}
\caption{מפת קווי השבר הייחודית לישראל לקראת קיר המסירות (\num{2026}--\num{2027}): מיקוד הכשל בשלוש חוליות הקצה ובלימת ההדבקה המערכתית על ידי חומת האש המוסדית \cite{boi2025stability,salelaw1974}}
\label{fig:israel_fault_line_map}
\end{figure}

\begin{enumerate}[label=\textbf{\arabic*.},rightmargin=0.6cm]
  \item \textbf{קו שבר א' – מלכודת משפרי הדיור והמשקיעים הממונפים (\en{Double Asset Lock}):} משפרי דיור שרכשו דירה חדשה במבצע \num{15/85} או בהלוואת בלון קבלן מתוך כוונה למכור את דירתם הקיימת סמוך למועד המסירה ב-\num{2026}--\num{2027}, מגלים כי שוק היד השנייה קפוא לחלוטין עקב התחרות מול מבצעי הקבלנים ומלאי ה-\num{85,900} דירות. ללא מכירת הדירה הישנה, הם חסומים רגולטורית מנטילת משכנתא מעל \pct{50}--\pct{70} וחשופים לקנס מס רכישה ולהפרת חוזה.
  \item \textbf{קו שבר ב' – יזמים קטנים ובינוניים (\en{Tier-2 / Tier-3}) הממונפים באשראי מזנין חוץ-בנקאי:} כדי להעמיד את ההון העצמי שנדרש על ידי הבנקים המלווים (\pct{20}--\pct{25} מעלות הפרויקט), יזמים רבים נטלו הלוואות השלמת הון (\en{Mezzanine}) מקרנות אשראי חוץ-בנקאיות בריבית של \pct{11}--\pct{14} (\en{Prime + 5\%--8\%}). בפרויקטים אלו, הבנק המלווה מובטח בשעבוד ראשון על הפרויקט, אך עיכוב של \num{6}--\num{12} חודשים בתקבולי המסירה או גל ביטולים של \pct{15}--\pct{20} מהחוזים מוחק לחלוטין את עודפי הפרויקט וגורם לחדלות פירעון של היזם ושל קרנות המזנין.
  \item \textbf{קו שבר ג' – משקי בית בעשירוני הביניים שנדחקים מעבר לתקרת ה-\en{PTI} במועד המסירה:} רוכשים שחתמו על עסקאות \num{20/80} ליניאריות ללא משכנתא מאושרת מראש ונדרשים לגייס במועד המסירה משכנתא של \num{1.8}--\num{2.4} מיליון ש"ח בריבית של \pct{5.5}, כאשר ההחזר החודשי חוצה את רף ה-\pct{40}--\pct{50} מהכנסתם הפנויה והבנק המסחרי מסרב להעמיד את ההלוואה.
\end{enumerate}

\begin{managerbox}[סיכום פרק ד' לקובעי המדיניות: הבחנה בין סיכון מערכתי לסיכון סקטוריאלי]
הניתוח ההשוואתי בפרק זה מוביל למסקנה אסטרטגית חדה: ישראל של \num{2026} אינה עומדת בפני \en{"}משבר ליהמן ברדרס\en{"} של קריסת בנקים מסחריים, בזכות חומת האש המשולשת של חוק המכר, הוראות ה-\en{LTV/PTI} של בנק ישראל ומשטר ה-\en{Full Recourse}. אולם עצם עוצמתה של חומת האש הבנקאית משמעותה ש\textbf{הבנקים מוגנים על ידי העברת מלוא הלם הנזילות אל החוליות החלשות בשרשרת}: רוכשי \en{"}על הנייר\en{"} שייאלצו למכור דירות במחירי מצוקה לפני טופס \num{4}, יזמי נדל"ן ממונפים שייקלעו לכינוס נכסים, וגופי אשראי חוץ-בנקאיים שהעמידו הלוואות מזנין ובלון ללא כריות ביטחון מספקות. פרק ה' ופרק ו' יכמתו בדיוק מתמטי את עוצמת הפגיעה בכל אחת מחוליות אלו.
\end{managerbox}
